\section{Bounded Literature Audit and Collision Analysis}

This section records the enlarged literature review used to stress-test the scope claim. It is a bounded technical audit, not a claim of an exhaustive systematic review. The accompanying \texttt{literature\_matrix.csv} contains 65 distinct scholarly records, a direct scholarly verification URL, a stable identifier, source tier, access date, role, boundary and reading depth. Fifteen works received a substantive full-paper or theorem/algorithm pass; the remaining 50 were checked bibliographically and against the passages needed for positioning. Three legitimate edition discrepancies are explicitly cross-checked rather than silently merged: ECP proceedings pagination, the 2006 workshop/2007 Springer chapter date, and one terminal-page discrepancy in an ECAI export. Every work below is cited in context, rather than inserted through an uncited wildcard command. The accompanying research note retains only the ten sources that carry its argument; this supplement carries the complete audit.

\subsection{Shortest Paths, Heuristic Search, and Planning Context}
The underlying distance computations descend from shortest-path dynamic programming and label-setting algorithms \citep{bellman1958,dijkstra1959}. Admissible best-first search and its standard textbook treatment establish the search-side interpretation of lower bounds \citep{hart1968,pearl1984}. Classical planning as heuristic search and the Fast Downward architecture place those ideas in the planning setting \citep{bonetgeffner2001,helmert2006}. State-equation, landmark, abstraction, and potential heuristics show that cost-optimal planning already contains several families of admissible lower bounds and optimization views \citep{bonet2013,karpasdomshlak2009,helmertdomshlak2009,seipppommerening2015}. These works motivate the setting, but they do not study how one externally changed action cost is transformed by a fixed sequence of saturation subtractions.

\subsection{Pattern Databases and Abstraction Heuristics}
Pattern databases originated as stored exact distances in an abstraction and were brought into optimal planning through automatic construction methods \citep{culberson1998,korf1997,edelkamp2001,edelkamp2006}. Additive and maximum combinations clarify when several abstractions may be combined without losing admissibility \citep{felner2004,holte2006,haslum2007,yang2008}. Optimal abstraction composition and planning-specific PDB selection then made cost allocation and pattern interaction explicit \citep{katzdomshlak2008,katzdomshlak2010,pommerening2013}.

Cartesian abstraction refinement supplies another route to exact abstract distances, including diverse additive collections and the later full treatment \citep{seipphelmert2013,seipphelmert2014,seipphelmert2018}. Merge-and-shrink develops a complementary abstraction family and analyzes label reduction and merge choices \citep{helmert2014,sievers2014,sievers2016}. Later work on refinement, factored tasks, transition generation, conditional effects, and abstraction merging broadens the construction techniques available to a planner \citep{speckseipp2022,abstractionfactored2024,cartesiantransitions2024,pozo2025,mergecartesian2025}. Planner ablations and the dissertation-level synthesis document how Cartesian abstractions and saturated partitions operate in complete systems \citep{seipp2024dissecting,seipp2018thesis}. Our constructions intentionally use the smallest possible two-state abstractions, so the result cannot be attributed to a difficult abstraction-building procedure.

\subsection{Cost Partitioning and Saturated Cost Partitioning}
General operator-cost partitioning, LP-based heuristics, state-dependent partitions, and Lagrangian formulations establish that allocating action cost among abstractions is an optimization problem with many admissible solutions \citep{fromnegative2015,lpbased2014,statedependent2016,lagrangian2019}. The development of saturated cost partitioning includes comparative evaluation, order optimization, tightening toward optimal partitions, pattern selection, subset saturation, and the complete journal treatment \citep{comparison2017,betterorders2017,narrowing2017,patternselection2019,subsetsaturated2019,scp2020}. We inherit the ordered recurrence and its admissibility argument from this line; neither is renamed as a contribution here.

Several nearby methods change what is optimized or represented. Online SCP changes which stored orders are evaluated; transition cost partitioning allocates at transition rather than label granularity; saturated post-hoc optimization and Dantzig--Wolfe decomposition optimize combinations rather than replaying one frozen ordered recurrence \citep{online2021,transitioncp2021,saturatedposthoc2021,dw2021}. Sensitivity analysis reuses LP solutions, probabilistic planning changes the task semantics, perfect SCP compresses order behavior, monotonic variants study adding or refining abstractions, and cost partitioning for sequence alignment transfers the technique to another domain \citep{sensitivity2023,probabilisticscp2023,sensitivity2024,perfect2025,monotonic2026,costpartitionmsa2024}. These are the strongest semantic neighbors. None, in the reviewed formulation, proves a bounded-incidence external-cost cascade for a fixed abstraction set and order or gives the table-bound exact-reuse rule analyzed here.

\subsection{Dynamic and Parametric Shortest Paths}
Incremental search has also been used inside Cartesian abstraction refinement \citep{incremental2020}. Lifelong Planning A*, D* Lite, focused D*, and later navigation/replanning analyses reuse shortest-path information as one graph changes \citep{lpa2004,dstar2002,stentz1995,koeniglichachev2005,hernandez2012}. General dynamic shortest-path algorithms maintain one graph under updates, while parametric shortest-path work studies structured variation of edge weights \citep{ramalingam1996,frigioni2000,demetrescu2004,karporlin1981,youngtarjanorlin1991}. Such an algorithm may accelerate a single rebuilt abstraction. It does not by itself decide which later abstraction receives which new residual costs, because those costs are outputs of earlier saturation stages. This distinction is why the relay remains a counterexample even when every individual PDB is trivial.

\subsection{Strongest Collisions and the Retained Boundary}
Table~\ref{tab:collision-audit} states the collision test used to keep the claim narrow. The negative-result route is important: incremental repair is not uniformly faster than recomputation, so the artifact reports relay cases where its exact certificate overhead loses rather than filtering them out \citep{hernandez2012}. Likewise, a familiar shortest-path optimality certificate would not be enough by itself. The retained contribution is the composition result: a sparse external update can travel and amplify through the residual pipeline, and a cached table must remain paired with the allocation that preserves it.

\begin{table*}[t]
\centering
\small
\begin{tabular}{@{}p{0.17\textwidth}p{0.34\textwidth}p{0.42\textwidth}@{}}
\toprule
Neighbor & Already established & Boundary retained here \\
\midrule
SCP and order methods \citep{comparison2017,scp2020,betterorders2017} & Ordered saturation, admissibility, and the importance of abstraction order. & External action costs change between searches while graphs and one order stay fixed; the question is propagation through that fixed recurrence. \\
Optimization and sensitivity \citep{dw2021,sensitivity2023,sensitivity2024} & Partition optimization and exact reuse regions for LP-based formulations. & No generic sensitivity or LP result is claimed; the repair result targets saved exact PDB values and their saturations. \\
Online and perfect SCP \citep{online2021,perfect2025} & Efficient evaluation or representation of many useful abstraction orders. & There is one frozen order; no maximum over orders or order-selection speedup is credited to the method. \\
Monotonic and transition variants \citep{transitioncp2021,monotonic2026} & Finer allocation units and behavior when abstraction collections change. & Labels remain shared across projected transitions and the abstraction collection is unchanged. \\
Dynamic shortest paths \citep{ramalingam1996,lpa2004,incremental2020} & Reuse of distances after changes in one graph. & Such reuse is an interchangeable within-table subroutine; it does not remove inter-table residual dependence. \\
This work & No new shortest-path algorithm, PDB construction, or cost-partition recurrence. & A tight degree-three Fibonacci family, a degree-two arbitrary relay, and exact ordered repair with explicit stale-value and zero-cycle controls. \\
\bottomrule
\end{tabular}
\caption{Collision audit. The right column is intentionally narrower than a priority claim: it states the model distinction supported by the proofs and artifact.}
\label{tab:collision-audit}
\end{table*}

The audit therefore supports a limited novelty statement only. It does not prove that no equivalent construction exists outside the reviewed set, and it does not turn the standard tight-edge condition into a new certificate theorem. It does show why the result is not subsumed merely by citing dynamic shortest paths, order selection, LP sensitivity, or abstraction refinement: all four omit the specific fixed-order residual coupling attacked by the constructions and controls.
